\section{Module \texttt{container}}

Le module \texttt{container} regroupe les conteneurs de mise en page : grille, boîte linéaire et pile de pages.

\subsection{Macros de Configuration}

\begin{docbox}{Macros \texttt{GRID\_CONFIG\_DEFAULT}, \texttt{BOX\_CONFIG\_DEFAULT}, \texttt{STACK\_CONFIG\_DEFAULT}}
\begin{lstlisting}[language=C]
#define GRID_CONFIG_DEFAULT ((grid_config){ \
    .column_spacing = 0, \
    .row_spacing = 0, \
    .homogeneous_columns = false, \
    .homogeneous_rows = false, \
    .style = WIDGET_STYLE_CONFIG_DEFAULT, \
})

#define BOX_CONFIG_DEFAULT ((box_config){ \
    .orientation = GTK_ORIENTATION_VERTICAL, \
    .spacing = 0, \
    .homogeneous = false, \
    .style = WIDGET_STYLE_CONFIG_DEFAULT, \
})

#define STACK_CONFIG_DEFAULT ((stack_config){ \
    .transition = GTK_STACK_TRANSITION_TYPE_NONE, \
    .duration_ms = 200, \
    .vhomogeneous = false, \
    .hhomogeneous = false, \
    .style = WIDGET_STYLE_CONFIG_DEFAULT, \
})
\end{lstlisting}
\end{docbox}

Ces macros garantissent des configurations minimales sûres et prévisibles. Elles évitent de laisser des paramètres de layout dans un état indéfini, ce qui simplifie le débogage de l'interface. Elles permettent aussi une initialisation rapide avant surcharge locale.

En pratique, ces valeurs par défaut évitent les erreurs de mise en page (espacements incohérents, transition non souhaitée, orientation oubliée). Elles sécurisent le démarrage de l'interface dès les premiers prototypes.

\subsection{Fonctions API}

\begin{lstlisting}[language=C]
GtkWidget *create_grid(const grid_config *config) {
    GtkWidget *grid = gtk_grid_new();
    if (config == NULL) return grid;

    gtk_grid_set_column_spacing(GTK_GRID(grid), config->column_spacing);
    gtk_grid_set_row_spacing(GTK_GRID(grid), config->row_spacing);
    gtk_grid_set_column_homogeneous(GTK_GRID(grid), config->homogeneous_columns);
    gtk_grid_set_row_homogeneous(GTK_GRID(grid), config->homogeneous_rows);
    apply_widget_style(grid, &config->style);
    return grid;
}

GtkWidget *create_stack(const stack_config *config) {
    GtkWidget *stack = gtk_stack_new();
    if (config == NULL) return stack;

    gtk_stack_set_transition_type(GTK_STACK(stack), config->transition);
    gtk_stack_set_transition_duration(GTK_STACK(stack), config->duration_ms);
    return stack;
}
\end{lstlisting}

La logique applique les propriétés seulement si \texttt{config} existe, ce qui protège le code appelant et facilite les usages simples.

Le module \texttt{container} joue un rôle clé dans l'ergonomie globale: il encapsule la disposition des widgets et isole les détails GTK liés au layout. Cette séparation améliore la lisibilité du code applicatif.

\newpage
